% Transtornos
    % sacar lo que pusimos en el cap 4 sobre clinica
    % modelo entrenaod con dep y ano para self-harm, se necesitarían trabajo interdicippliario o a futuro para explorar más la relación entre estas problematicas porque los resultados experimentales sugieren que es la relación es lo suficientemetne fuerte como para poder detectar selfharm.

% resumen
En este trabajo hemos introducido, desarrollado, evaluado y analizado a SS3, un clasificador de texto diseñado con el objetivo de abordar escenarios de detección anticipada de riesgos en entornos dinámicos, como el de las redes sociales, de forma eficiente y unificada.
En particular, vimos que este tipo de escenarios es especialmente desafiante para los modelos convencionales debido a que abordarlas involucra, al menos, tres aspectos clave: \emph{clasificación de secuencias}, \emph{clasificación anticipada} y \emph{transparencia/explicabilidad}.
Por consiguiente, SS3 fue diseñado para integrar, respectivamente, las siguientes tres capacidades: \emph{trabajo/funcionamiento incremental}, \emph{soporte de mecanismos de toma de decisión anticipada}, y \emph{capacidad de explicar visualmente las razones detrás de sus decisiones}.

% evaluación y resultados
La evaluación del modelo se llevó a cabo, primeramente, en la tarea de detección anticipada de depresión de la primer edición del eRisk (2017).
Luego, se lo evaluó junto a $\tau$-SS3, una extensión del modelo original, en las dos tareas del eRisk 2018, detección anticipada de anorexia y detección anticipada de depresión. Finalmente, ambos modelos fueron puestos a prueba participando en las dos tareas del eRisk 2019, detección anticipada de anorexia y detección anticipada de conductas autodestructivas.
Los resultados obtenidos fueron, en relación a los otros modelos, consistentemente buenos a lo largo de todas las tareas.
En particular, el modelo obtuvo los mejores valores $ERDE_5$ y $ERDE_{50}$ en las distintas tareas abordadas y, en relación a las nuevas medidas introducidas en el eRisk 2019, obtuvo el mejor valor para la medida $F_{latency}$ en la detección anticipada de conductas autodestructivas, y el mejor desempeño general para las medidas basadas en ranking en ambas tareas de esa edición del eRisk ---habiendo conseguido los mejores valores $P@10$ y  $NDCG@10$, consistentemente, en los cuatros rankings y en las dos tareas utilizadas para evaluarlo, junto a los mejores valores $NDCG@100$ para los rankings generados con 1 y 100 publicaciones, y los segundos mejores para los generados con 500 y 1000, en ambas tareas.
Además, como SS3 fue diseñado para poder trabajar incrementalmente, permitiéndole procesar secuencias en tiempo $O(n)$ con respecto a la longitud de las mismas, los resultados también mostraron que el modelo es computacionalmente eficiente al haber sido el más rápido en las distintas tareas.
Finalmente, los resultados sugieren que SS3 puede llegar a ser un modelo robusto, ya que, independientemente de los datos utilizados para entrenarlo, su desempeño se mantuvo consistente a lo largo de las distintas tareas y medidas de desempeño utilizadas para evaluarlo, a pesar de haber utilizado el mismo valor de hiperparámetros para todas ellas.
Además, otro aspecto interesante a destacar es su simpleza e independencia del dominio, ya que no utilizó ningún tipo de representaciones de documentos compleja, o manualmente creada, ni mecanismos complejos de clasificación anticipada específicos para cada dominio abordado ---lo que contrasta con otros modelos que requerirían de un proceso costoso para adaptarlos a nuevos problemas.

% trabajo a futuro
% Sin embargo, si bien SS3 tuvo un desempeño superior en relación a los otros modelos, su desempeño se encontró lejos de ser perfecto.
% Por ejemplo, si bien los valores $ERDE$ fueron los menores de entre los demás modelos, no fueron lo suficientemente cercanos a $0$ como para indicar un desempeño sobresaliente en términos absolutos ---valores cercanos a $0$ indicarían una clasificación rápida y precisa, con un porcentaje de error cercano a nulo.
% Lo mismo se pudo apreciar en relación a la medida $F_{latency}$, en donde el valor obtenido para la tarea de detección anticipada de anorexia del eRisk 2019 no fue el mejor, o en la tarea de detección de conductas autodestructivas, en donde si bien obtuvo el mejor valor ($0.52$), dicho valor se encontró lejos de ser sobresaliente, es decir, cercano a $1$.
% Este hecho nos señala que, si bien los resultados obtenidos fueron alentadores, aún se requiere de mucho trabajo para que, de hecho, se puedan desarrollar sistemas realmente eficientes en detectar automática y anticipadamente situaciones de riesgo, en la práctica.
Sin embargo, si bien los resultados obtenidos fueron alentadores en relación a los otros modelos, todavía se tiene mucho margen de mejora en términos absolutos como para que, de hecho, se puedan desarrollar sistemas realmente eficientes en detectar automática y anticipadamente situaciones de riesgo, en la práctica.
Por este motivo, como trabajo futuro, planeamos explorar diferentes aspectos vinculados a mejorar el desempeño de clasificación anticipada del modelo.
En este contexto, creemos que los resultados obtenidos en las medidas basadas en ranking son claves, ya que nos señalan que SS3, de hecho, aprendió a valorar correctamente el nivel de riesgo de los usuarios y que, por lo tanto, si no obtuvo un mejor desempeño en la clasificación no se debió a una mala estimación de dicho valor, sino a la política utilizada para decidir \emph{cuándo} clasificarlos anticipadamente en base al mismo.
En particular, la política utilizada mostró ser ``demasiado precoz'' al  priorizar, con diferencia, velocidad por sobre exactitud ---habiendo clasificado a los posibles usuarios en riesgo luego de procesar sólo 3 de sus publicaciones, en promedio.
En consecuencia, como trabajo futuro exploraremos mejores políticas de clasificación anticipada, por ejemplo, creemos que retrasar la clasificación hasta que el nivel de riesgo estimado del usuario sea lo suficientemente alto o utilizar una política que haga uso de información global, relacionada al nivel de riesgo de todos los usuarios, podría ayudar a mejorar significativamente el rendimiento del clasificador.
Además, buscaremos distintos valores de hiperparámetros que nos permitan mejorar el rendimiento en términos de la nueva medida $F_{latency}$, ya que los utilizados fueron seleccionados en relación a la medida $ERDE$. %---la cual prioriza, en gran medida, velocidad y \emph{recall} por sobre precisión.
Asimismo, también realizaremos pruebas utilizando diferentes \emph{operadores de resumen}, $\oplus_j$, ya que en este trabajo no se lo consideró y podría tener un impacto positivo en el desempeño del modelo. Estos \emph{operadores de resumen} podrían hasta incluso llegar a ser funciones complejas que sean capaces de realizar operaciones más avanzadas que una simple operación aritmética dada, ya que en principio se podría utilizar cualquier algoritmo que tome como entrada una lista de vectores y retorne un único vector, por ejemplo, estos \emph{operadores de resumen} se podrían hasta incluso aprender utilizando modelos de aprendizaje automático.
Finalmente, al igual que con $\tau$-SS3, también exploraremos otras maneras de extender el modelo, por ejemplo, se podrían explorar formas alternativas para calcular el \emph{valor global}, $gv(w,c)$, ya sea utilizando nuevas ecuaciones o realizando mejoras sobre las actuales.
Además, en relación a una nueva extensión, creemos que utilizar información sobre la coocurrencia de las palabras podría contribuir positivamente en el desempeño del modelo para tareas como las abordadas, ya que las palabras que coocurran con palabras tales como ``I'', ``myself'' o ``my'' podrían ser muy discriminativas, permitiéndole al modelo detectar oraciones completas en las que el autor se refiere a sí mismo ---esta extensión se podría implementar, por ejemplo, utilizando un grafo de coocurrencia por cada clase, en donde cada arista de los grafos esté ponderada por su \emph{valor global}, aprendido en base al \emph{valor local} de dicha artista a lo largo de los distintos grafos.

Por otro lado, ya que el modelo introducido es en esencia un clasificador de texto de propósito general, este trabajo doctoral habilita nuevas líneas de investigación no vinculadas específicamente a la detección anticipada de riesgos, tales como las vinculadas a examinar distintas propiedades del clasificador, su desempeño en diferentes tipos de tareas, o su integración con otros modelos.
Por ejemplo, en este trabajo pudimos observar, tanto cualitativamente por medio de las nubes de palabras como cuantitativamente en base a los resultados basados en ranking, que el \emph{valor global} puede ser un buen estimador de la importancia de las palabras en relación a cada clase, por lo que creemos que podría jugar un papel importante, por ejemplo, como método de pesado de término o como método de selección de características. Por lo tanto, como trabajo futuro utilizaremos al \emph{valor global} como método de pesado de términos y selección de características, comparándolo con otros enfoques bien conocidos, como la \emph{ganancia de información}, \emph{chi-cuadrado} ($\chi^2$), \emph{tf-idf}, entre otros.
De igual modo, examinaremos la integración de los \emph{vectores de confianza} producidos por SS3 en otros modelos, por ejemplo, al utilizarlos para construir nuevas representaciones vectoriales, densas, de los documentos de entrada, o al utilizarlos como ``embeddings supervisados'' de distintos bloques de texto\footnote{Por ejemplo, embeddings de palabras, embeddings de oraciones, embeddings de documentos, etc.} ---los cuales luego podrían ser integrados, por ejemplo, a modelos secuenciales tales como redes neuronales convolucionales o distintos tipos de redes neuronales recurrentes, como las LSTM o las GRU.
Además, sería interesante probar la robustez del modelo ante ciertos problemas bien conocidos dentro del área.
Por ejemplo, en base a que el modelo ve a cada clase como una ``gran bolsa de palabras'' en donde el número de documentos utilizados para construirlas durante el entrenamiento es irrelevante,\footnote{Siempre que la frecuencia relativa de las palabras sea la adecuada.} y en base a los resultados obtenidos, en donde los conjuntos de datos utilizados fueron altamente desbalanceados, creemos que el modelo podría tener cierta robustez ``natural'' ante el conocido \emph{problema de desbalance de clases}, por lo que realizaremos pruebas en distintas tareas para poder cuantificar dicha robustez.
Del mismo modo, al aprender a ponderar las palabras globalmente, sin vincularse directamente a cada documento, sus propiedades o su origen, el modelo tal vez podría llegar a tener cierta robustez para lidiar con el problema de la \emph{clasificación en dominios cruzados}, como podrían sugerirlo los resultados obtenidos por el modelo \emph{UNSL\#3} en la tarea de detección de conductas autodestructivas del eRisk 2019, el cual mostró un buen desempeño en relación a los otros modelos, a pesar de haber sido entrenado con datos de anorexia y depresión.
Finalmente, sería interesante medir el desempeño del clasificador en distintas tareas de clasificación de documentos tradicionales, no sólo en clasificación de etiquetado único sino también de multi-etiquetado,\footnote{En inglés, ``single-label clasification'' y ``multi-label clasification'', respectivamente.} ya que no hay ninguna restricción que le impida al modelo asignarle más de una clase al documento de entrada, por lo que también examinaremos el uso de distintas políticas para seleccionar múltiples clases a partir del \emph{vector de confianza} del documento.

\

Por último, dada la naturaleza de la problemática abordada, creemos adecuado concluir este trabajo, al igual que en el \autoref{cap:ss3}, volviendo a remarcar algunos aspectos vinculados a las implicaciones y consideraciones clínicas, y éticas, en relación a su aplicación en la vida real.
% Por último, dada la naturaleza de la problemática abordada, creemos adecuado concluir este trabajo, al igual que en el \autoref{cap:ss3}, con algunos aspectos vinculados a las implicaciones y consideraciones clínicas y éticas del trabajo realizado, en relación a su aplicación en la vida real.
% concluir con algunos aspectos vinculados a posibles aplicaciones las implicaciones y consideraciones clínicas del trabajo realizado.
% mencionar nuevamente, como en la \autoref{sec:implications_clinical}, algunos aspectos vinculados a las implicaciones y consideraciones clínicas del trabajo realizado.
% Los métodos de detección automática pueden ayudar a identificar a personas en riesgo a través del monitoreo pasivo, a gran escala, de las redes sociales, y en el futuro pueden complementar los procedimientos de detección existentes~\cite{guntuku2017detecting}.
% En este contexto, el modelo presentado en este trabajo se pensó como una herramienta que podría ser utilizada, en un futuro, para desarrollar sistemas de monitoreo pasivo a gran escala con los que ayudar a detectar usuarios en riesgo analizando sus patrones lingüísticos. Por ejemplo, dicho sistema podría filtrar usuarios presentando posibles candidatos a profesionales de la salud mental para ser analizados manualmente ---tal vez por medio de un ranking de candidatos ordenados por nivel de riesgo, junto con información visual e interactiva.
El modelo presentado en este trabajo se pensó como un marco de trabajo que podría ser utilizada, a futuro, para desarrollar sistemas de monitoreo pasivo a gran escala con los que ayudar a detectar usuarios en riesgo analizando sus patrones lingüísticos. Por ejemplo, dicho sistema podría filtrar usuarios presentando posibles candidatos a profesionales de la salud mental para ser analizados manualmente ---tal vez por medio de un ranking de candidatos ordenados por nivel de riesgo, junto con información visual e interactiva.
El aspecto de ``monitoreo pasivo a gran escala'' podría estar respaldado por la naturaleza paralelizable e incremental de SS3,\footnote{Sólo sería necesario almacenar un pequeño vector bidimensional por cada usuario, su \emph{vector de confianza}, el cual se actualizaría poco a poco a medida que se genere nuevo contenido.} mientras que la presentación de ``información visual e interactiva'' a los profesionales, por su capacidad de explicar visualmente las razones detrás de sus decisiones.
Está claro que el trabajo presentado en esta tesis doctoral \emph{no persiguió un objetivo de diagnóstico automático}, sino el desarrollo, a futuro, de herramientas de monitoreo complementarias a otros métodos bien establecidos dentro del estudio de la salud mental.
No obstante, los problemas éticos y legales sobre la privacidad y protección de datos que involucraría llevar realmente a la práctica sistemas de monitoreo de este tipo siguen siendo problemas abiertos de investigación~\cite{guntuku2017detecting}. \hfill$\square$

% Finalmente, aunque la interpretación clínica de los resultados no se abordó en nuestro trabajo, es importante resaltar algunos puntos importantes.
% En primer lugar, fue interesante observar que la mayoría de las 100 palabras identificadas por SS3, seleccionadas por mayor \emph{valor global}, se corresponden perfectamente a las categorías habitualmente identificadas en otros estudios más clínicos sobre depresión ---tales como palabras relacionadas a síntomas, tratamientos, divulgación/revelación (disclosure), relaciones/vida (relationships-life) \cite{ de2013social}.

